\chapter*{แผนการบริหารการสอนและประมวลรายวิชา}
\addcontentsline{toc}{chapter}{แผนการบริหารการสอนและประมวลรายวิชา}

เอกสารคำสอนและคู่มือวิจัยฟิสิกส์ระดับมหาวิทยาลัย ออกแบบขึ้นเพื่อใช้ในการเรียนการสอนวิชาฟิสิกส์ขั้นกลางถึงขั้นสูงสำหรับนักศึกษาระดับปริญญาตรี โดยมีแผนการสอนและการประเมินผลสัมฤทธิ์ตามกรอบมาตรฐานคุณวุฒิระดับอุดมศึกษา ดังแสดงในตารางต่อไปนี้

\vspace{0.4cm}
\begin{table}[htbp]
\centering
\small
\begin{tabularx}{\textwidth}{>{\centering\arraybackslash}p{1.3cm} >{\raggedright\arraybackslash}p{3.7cm} >{\raggedright\arraybackslash}X >{\raggedright\arraybackslash}p{2.6cm}}
\toprule
\textbf{สัปดาห์} & \textbf{ชื่อหัวข้อบทเรียน} & \textbf{สาระการเรียนรู้และกิจกรรมวิจัย} & \textbf{การประเมินผล} \\
\midrule
1--2 & หน่วย มาตรฐาน และการวิเคราะห์ ความไม่แน่นอน & ระบบหน่วยเอสไอใหม่ (CODATA), การกระจายตัวแบบเกาส์เซียน และประเมินความไม่แน่นอนตาม GUM & โจทย์ปัญหา C-C-A-F \\
3--4 & พลศาสตร์ขั้นสูง และการเคลื่อนที่ แบบไม่เชิงเส้น & แรงต้านอากาศแปรผันตามอัตราเร็ว ($v$ และ $v^2$), ความเร็วปลาย และจำลองวิถีโค้งด้วย Runge-Kutta & การจำลองด้วย Python \\
5--6 & ทฤษฎีศักย์ งาน พลังงาน และเสถียรภาพ & สนามแรงอนุรักษ์, ฟังก์ชันพลังงานศักย์, สมดุลเสถียร/ไม่เสถียร และแผนภาพเฟส (Phase Space) & โจทย์ปัญหา P-S-C-T \\
7--8 & ทฤษฎีสนามแม่เหล็กไฟฟ้า และสมการแมกซ์เวลล์ & กฎของเกาส์ ฟาราเดย์ แอมแปร์-แมกซ์เวลล์ รูปอนุพันธ์ย่อย และสมการคลื่นแม่เหล็กไฟฟ้าในสุญญากาศ & สอบวัดผลกลางภาค \\
9--10 & ทัศนศาสตร์เชิงคลื่น และการแทรกสอดเลเซอร์ & คลื่นแสงอาพันธ์, การแทรกสอดแบบยังก์สองช่อง, การเลี้ยวเบนฟรอนโฮเฟอร์ และโครงสร้างผลึก & รายงานการทดลอง \\
11--12 & กลศาสตร์ควอนตัม และแบบจำลองอะตอม & ทวิภาวะคลื่น-อนุภาค, สมการชเรอดิงเงอร์, อนุภาคในกล่องศักย์ และการทะลุผ่านกำแพงศักย์ & การวิเคราะห์เชิงตัวเลข \\
13--14 & ฟิสิกส์ของแข็ง และวัสดุแม่เหล็กขั้นสูง & โครงสร้างผลึก, ทฤษฎีแถบพลังงาน, อันตรกิริยาแลกเปลี่ยน $J$ และการคำนวณ $\text{LSDA}+U$ ใน $\text{MnO}$ & สัมมนางานวิจัย \\
15 & ระเบียบวิธีวิจัย และฟิสิกส์เชิงคำนวณ & การแก้สมการเชิงอนุพันธ์ย่อย, Physics-Informed Neural Networks (PINNs) และการนำเสนอบทความวิจัย & โครงงานวิจัยเดี่ยว \\
\bottomrule
\end{tabularx}
\caption*{ตารางแสดงแผนการจัดการเรียนรู้ 15 สัปดาห์}
\end{table}
